% Section: entailment

\documentclass[../../../include/open-logic-section]{subfiles}

\begin{document}

\olfileid[fr]{nml}{syn}{ent}

\olsection{Conséquence logique}

\begin{explain}
  La notion de vérité en un monde permet de définir une relation
  de conséquence logique entre !!{formula}s. Une !!{formula}~$!B$
  entraîne~$!A$ si et seulement si, chaque fois que $!B$ est vraie,
  $!A$ l'est aussi. « Chaque fois » signifie ici : quel que soit
  le modèle considéré et quel que soit le monde considéré dans ce
  modèle.
\end{explain}

\begin{defn}
  Si $\Gamma$ est un ensemble de !!{formula}s et $!A$ une
  !!{formula}, alors $\Gamma$ \emph{entraîne}~$!A$, ce que nous
  notons $\Gamma \Entails !A$, si et seulement si, pour tout
  modèle $\mModel{M} = \tuple{W, R, V}$ et tout monde $w \in W$,
  dès que $\mSat{M}{!B}[w]$ pour chaque $!B \in \Gamma$, on a
  $\mSat{M}{!A}[w]$. Si $\Gamma$ ne contient qu'une
  !!{formula}~$!B$, nous écrivons $!B \Entails !A$.
\end{defn}

\begin{ex}
  Pour montrer qu'une !!{formula} en entraîne une autre, il faut
  raisonner sur tous les modèles à partir de la définition de
  $\mSat{M}{}[w]$. Montrons par exemple que
  $p \lif \Diamond p \Entails \Box\lnot p \lif \lnot p$.
  Soient un modèle $\mModel{M} = \tuple{W, R,
    V}$ et un monde $w \in W$, et supposons
  $\mSat{M}{p \lif \Diamond p}[w]$. Il faut établir
  $\mSat{M}{\Box\lnot p \lif \lnot p}[w]$. Supposons le contraire.
  Alors $\mSat{M}{\Box\lnot p}[w]$ et
  $\mSat/{M}{\lnot p}[w]$. Puisque
  $\mSat/{M}{\lnot p}[w]$, nous déduisons
  $\mSat{M}{p}[w]$. Par hypothèse,
  $\mSat{M}{p \lif \Diamond p}[w]$, donc
  $\mSat{M}{\Diamond p}[w]$. Par définition de
  $\mSat{M}{\Diamond p}[w]$, il existe un monde $w'$ tel que
  $Rww'$ et $\mSat{M}{p}[w']$. Or
  $\mSat{M}{\Box \lnot p}[w]$ donne aussi
  $\mSat{M}{\lnot p}[w']$, contradiction.

  Pour montrer qu'une !!{formula}~$!B$ n'en entraîne pas une
  autre~$!A$, il faut un contre-exemple : un modèle
  $\mModel{M} = \tuple{W, R, V}$ et un monde~$w \in W$ tels que
  $\mSat{M}{!B}[w]$, mais $\mSat/{M}{!A}[w]$.
  Montrons ainsi que $p \lif \Diamond p \Entails/
  \Box p \lif p$. Considérons le modèle de la
  \olref{fig:counterex}. Nous avons
  $\mSat{M}{\Diamond p}[w_1]$, donc
  $\mSat{M}{p \lif \Diamond p}[w_1]$. Toutefois,
  $\mSat{M}{\Box p}[w_1]$ et $\mSat/{M}{p}[w_1]$;
  par conséquent, $\mSat/{M}{\Box p \lif p}[w_1]$.
  \begin{figure}
  \begin{center}
    \begin{tikzpicture}[modal]
      \node[world] (w1) [label=right:\mFalse{p}]{$w_1$}; 
      \node[world] (w2) [label=right:\mTrue{p}, above left=of w1]{$w_2$}; 
      \node[world] (w3) [label=right:\mTrue{p}, above right=of w1] {$w_3$};
      \draw[->] (w1) to (w2);
      \draw[->] (w1) to (w3);
    \end{tikzpicture}
  \end{center}
  \caption{Contre-exemple à $p \lif \Diamond p
  \Entails/ \Box p \lif p$.}
  \ollabel{fig:counterex}
  \end{figure}
  
  Des contre-exemples très simples suffisent souvent. Le modèle
  $\mModel{M'} = \tuple{W', R', V'}$ avec $W' = \{w\}$,
  $R' = \emptyset$ et $V'(p) = \emptyset$ en fournit un autre.
  En effet, comme $\mSat/{M'}{p}[w]$,
  $\mSat{M'}{p \lif \Diamond p}[w]$. Aucun monde n'étant
  accessible depuis~$w$, nous avons $\mSat{M'}{\Box p}[w]$,
  donc $\mSat/{M'}{\Box p \lif p}[w]$.
\end{ex}

\begin{prob}
  Montrez que $\Box (!A \land !B) \Entails \Box !A$.
\end{prob}

\begin{prob}
  Montrez que $\Box (p \lif q) \Entails/ p \lif \Box q$ et que
  $p \lif \Box q \Entails/\Box (p \lif q)$.
\end{prob}
\begin{editorial}La légende du contre-exemple affichait dans la source le symbole de conséquence logique ; elle porte ici celui de non-conséquence, conformément au texte et au modèle. Le modèle à un seul monde était écrit comme un ensemble de trois composantes ; il est écrit ici comme le triplet ordonné prévu par la définition.\end{editorial}
\end{document}
